\chapter{Problem Formulation}
\label{ch:formulation}

This chapter sets up the optimization problem eALS solves, starting from
the plain matrix-factorization model and Hu et al.'s whole-data objective,
before introducing the popularity-aware weighting scheme that
distinguishes \citet{he2016eals} from earlier work and motivates the fast
solver derived in Chapter~\ref{ch:methods}.

\section{Matrix factorization for implicit feedback}

Let $R \in \{0,1\}^{M\times N}$ be the user-item interaction matrix, with
$M$ users and $N$ items. $\mathcal{R}$ denotes the set of user-item pairs
with $r_{ui}=1$ (an observed, positive interaction); every other entry is
zero. This is the crux of the implicit-feedback setting: a zero does not
mean ``the user dislikes the item'', it means ``no signal was observed''.
Matrix factorization approximates each entry as an inner
product of a user latent vector $p_u \in \mathbb{R}^K$ and an item latent
vector $q_i \in \mathbb{R}^K$, where $K$ is the number of latent factors, a
hyperparameter fixed before training and shared by every user and item:
\[
  \hat{r}_{ui} = \langle p_u, q_i \rangle = p_u^\top q_i .
\]
Stacking the vectors gives $P \in \mathbb{R}^{M\times K}$ and
$Q \in \mathbb{R}^{N\times K}$, the two matrices this project actually
learns.

\medskip
\citet{hu2008implicit} formulate learning as a weighted regression over
\emph{all} $M\times N$ entries, observed and missing alike:
\[
  J = \sum_{u=1}^{M}\sum_{i=1}^{N} w_{ui}(r_{ui}-\hat r_{ui})^2
      + \lambda\Big(\sum_u \lVert p_u\rVert^2 + \sum_i \lVert q_i\rVert^2\Big),
\]
where $w_{ui}$ is a confidence weight and $\lambda$ is the usual $L_2$
regularization strength. Whole-data learning of this kind is more faithful
than sampling a handful of negative examples per user, but it is also the
source of the whole efficiency problem this report is about: naively, every
gradient step touches all $MN$ entries, most of which are the missing
zeros. Hu et al.\ make this tractable by setting $w_{ui}$ to a single
constant $w_0$ for every missing pair, which allows the missing-data
contribution to be pre-computed once via $Q^\top Q$ and reused for every
user.

\section{Popularity-aware weighting}
\label{sec:popweight}

\citet{he2016eals} argue that this uniform-weight assumption discards
information most content providers already have: items that are frequently
recommended but rarely interacted with are more likely to be genuine
negatives than obscure items nobody has seen. They replace the single
missing-data term with a per-item confidence $c_i$, giving the objective
actually implemented in this project:
\begin{equation}
  \label{eq:objective}
  L = \sum_{(u,i)\in\mathcal{R}} w_{ui}(r_{ui}-\hat r_{ui})^2
      + \sum_{u=1}^{M}\sum_{i \notin \mathcal{R}_u} c_i\, \hat r_{ui}^{\,2}
      + \lambda\Big(\sum_u \lVert p_u\rVert^2 + \sum_i \lVert q_i\rVert^2\Big),
\end{equation}
where $\mathcal{R}_u$ is the set of items $u$ has interacted with, and,
symmetrically, $\mathcal{R}_i$ denotes the set of users who have
interacted with item $i$; both are used repeatedly in the rest of this
report and in Chapter~\ref{ch:methods} in particular. The per-item
confidence is derived from item popularity $f_i = |\mathcal{R}_i|
/ \sum_j |\mathcal{R}_j|$ (item $i$'s share of all interactions in the
training set):
\begin{equation}
  \label{eq:ci}
  c_i = c_0\, \frac{f_i^{\alpha}}{\sum_{j} f_j^{\alpha}},
\end{equation}
with two hyperparameters: $c_0$ sets the overall strength of the
missing-data penalty, and $\alpha$ controls how sharply popularity is
rewarded. Setting $\alpha=0$ collapses \eqref{eq:ci} back to Hu et al.'s
uniform weighting ($c_i \equiv c_0/N$ for every item), which makes $\alpha$
a convenient dial between the two papers' assumptions and is exploited
directly in this project's ablation study (Section~\ref{sec:weighting}).

\section{Why the naive solver does not scale}

Minimizing \eqref{eq:objective} with respect to a single user vector $p_u$,
holding everything else fixed, has a closed-form ridge-regression solution
that requires inverting a $K \times K$ matrix. Doing this once per user and
once per item costs $O((M+N)K^3)$ per pass over the data, on top of an
$O(|\mathcal{R}|K^2)$ term for the observed entries. The $K^3$ term
dominates as soon as a moderately large number of latent factors is used,
which is exactly the regime that gives MF its representational power.
Hu et al.'s uniform-weight trick removes an $O(N)$ factor from the
missing-data term but cannot remove the $K^3$ matrix inversion, because
that cost comes from updating a whole vector at once, independently of how
the missing data is weighted. Section~\ref{sec:ealsdesign} of the next
chapter derives the reformulation that eALS uses to eliminate this term
entirely, which is the paper's central algorithmic contribution and the
piece of software this report builds and evaluates.
